\section{Model-wise Double Descent}
\label{sec:model-dd}

\begin{figure}[h]
    \centering
    \begin{subfigure}[t]{.47\textwidth}
        \centering
          \includegraphics[width=\textwidth]{cifar100_resnet_pVar}
          \caption[width=0.8\linewidth]{{\bf CIFAR-100.}
          There is a peak in test error even with no label noise.
          } \label{fig:resnet-cifar-left}
    \end{subfigure}\hfill
    \begin{subfigure}[t]{.47\textwidth}
        \centering
          \includegraphics[width=\textwidth]{cifar10_resnet_pVar}
          \caption[width=0.8\linewidth]{{\bf CIFAR-10.}
        There is a ``plateau'' in test error around the interpolation point with no label noise,
        which develops into a peak for added label noise.
          }
    \end{subfigure}
    \caption{{\bf Model-wise double descent for
    ResNet18s.} Trained on CIFAR-100 and CIFAR-10, with varying label noise. Optimized using Adam with LR $0.0001$ for 4K epochs, and data-augmentation.
    }
    \label{fig:resnet-cifar}
\end{figure}

In this section, we study the test error of models of increasing size,
when training to completion (for a fixed large number of optimization steps). 
We demonstrate model-wise double descent across different architectures,
datasets, optimizers, and training procedures.
The critical region exhibits distinctly different test behavior
around the interpolation point and there is often a peak in test error that becomes more prominent in settings with label noise.

For the experiments in this section (Figures~\ref{fig:resnet-cifar}, \ref{fig:cnn-cifar},
\ref{fig:more-mdd1appendix},
\ref{fig:mcnn_cf100},
\ref{fig:more-mdd2}), notice that all modifications which increase the interpolation threshold
(such as adding label noise, using data augmentation, and increasing the number of train samples)
also correspondingly shift the peak in test error towards larger models.
Additional plots showing the early-stopping behavior of these models,
and additional experiments showing double descent
in settings with no label noise (e.g. Figure~\ref{fig:app-cifar100-clean})
are in Appendix~\ref{subsec:appendix_model}.
We also observed model-wise double descent for adversarial training,
with a prominent robust test error peak even in settings without label noise.
See Figure~\ref{fig:adv_training} in Appendix~\ref{subsec:appendix_model}.

\begin{figure}[h]
    \centering
    \begin{subfigure}{.47\textwidth}
        \centering
  \includegraphics[width=\textwidth]{cifar10_mcnn_noaug_pVar}
        \caption{Without data augmentation.}
    \end{subfigure}\hfill
    \begin{subfigure}{.47\textwidth}
        \centering
  \includegraphics[width=\textwidth]{cifar10_mcnn_aug_pVar_trunc}
        \caption{With data augmentation.}
    \end{subfigure}
    \caption{{\bf Effect of Data Augmentation.}
    5-layer CNNs on CIFAR10, with and without data-augmentation.
    Data-augmentation shifts the interpolation threshold to the right,
    shifting the test error peak accordingly.
    \SGDTxt~
     See Figure~\ref{fig:mcnn_aug_large} for larger models.
    }\label{fig:cnn-cifar}
\end{figure}

\begin{figure}[h]
\centering
\captionsetup{calcwidth=.9\linewidth}
\begin{minipage}[t]{0.47\textwidth}
  \centering
  \includegraphics[width=\textwidth]{cifar10_mcnn_sgd_adam}
  \caption[width=0.8\linewidth]{{\bf SGD vs. Adam.}
  5-Layer CNNs on CIFAR-10 with no label noise, and no data augmentation.
  Optimized using SGD for 500K gradient steps, and Adam for 4K epochs.}\label{fig:more-mdd1appendix}
\end{minipage}\hfill
\begin{minipage}[t]{.47\textwidth}
  \centering
  \includegraphics[width=\textwidth]{cifar100_mcnn_noaug}
    \caption{{\bf Noiseless settings.} 
    5-layer CNNs on CIFAR-100 with no label noise;
    note the peak in test error.
    Trained with SGD and no data augmentation.
    See Figure~\ref{fig:app-cifar100-clean-sgd}
    for the early-stopping behavior of these models.
    }
    \label{fig:mcnn_cf100}
\end{minipage}
\end{figure}

\begin{figure}[h]
    \centering
    \begin{minipage}[c]{0.5\textwidth}
        \includegraphics[width=\textwidth]{nlp_de_en_fr}
    \end{minipage}\hfill
    \begin{minipage}[c]{0.45\textwidth}
      \caption[0.9\textwidth]{{\bf Transformers on language translation tasks:} Multi-head-attention encoder-decoder Transformer model trained for 80k gradient steps with labeled smoothed cross-entropy loss on \iwslt\ German-to-English (160K sentences) and WMT`14 English-to-French (subsampled to 200K sentences) dataset. Test loss is measured as per-token perplexity.}
      \label{fig:more-mdd2}
    \end{minipage}
\end{figure}


% \begin{figure}
% \centering
% \begin{minipage}[t]{.5\textwidth}
% \captionsetup{width=.9\linewidth}
%   \centering
%   \includegraphics[width=\textwidth]{nlp_de_en_fr}
%       \caption{{\bf Transformers on language translation tasks:} Multi-head-attention encoder-decoder Transformer model trained for 80k gradient steps with labeled smoothed cross-entropy loss on \iwslt\ German-to-English (160K sentences) and WMT`14 English-to-French (subsampled to 200K sentences) dataset. Test loss is measured as per-token perplexity.}
%       \label{fig:more-mdd2}
% \end{minipage}%
% \begin{minipage}[t]{.5\textwidth}
% \captionsetup{width=.9\linewidth}
%   \centering
%   \includegraphics[width=\textwidth]{cifar100_mcnn_noaug}
%     \caption{Model-wise double descent for 5-layer CNNS on CIFAR-100,
%     and no label noise.
%     Note the peak in test error occurs even with no label noise.
%     Trained with SGD for 1e6 steps, and no data augmentation.
%     See Figure~\ref{fig:app-cifar100-clean-sgd}
%     for the early-stopping behavior of these models.
%     }
%     \label{fig:mcnn_cf100}
% \end{minipage}
% \end{figure}






\paragraph{Discussion.}
Fully understanding the mechanisms behind model-wise double descent
in deep neural networks remains an important open question.
However, an analog of model-wise double descent occurs even
for linear models. A recent stream of theoretical works
analyzes this setting
(\cite{bartlett2019benign,muthukumar2019harmless,belkin2019two,mei2019generalization,hastie2019surprises}).
We believe similar mechanisms may be at work in deep neural networks.

Informally, our intuition is that for model-sizes at the interpolation threshold, there is effectively
only one model that fits the train data and this interpolating model is very sensitive
to noise in the train set and/or model mis-specification.
%That is, these models have ``high variance'' in the classical sense.
That is, since the model is just barely able to fit the train data,
forcing it to fit even slightly-noisy or mis-specified
labels will destroy its global structure, and result in high test error.
(See Figure~\ref{fig:ens_resnet} in the Appendix for an experiment
demonstrating this noise sensitivity, by showing that ensembling helps significantly
in the critically-parameterized regime).
However for over-parameterized models,
there are many interpolating models that fit the train set,
and SGD is able to find one that ``memorizes''
(or ``absorbs'') the noise while still performing well on the distribution.

The above intuition is theoretically justified for linear models.
In general, this situation manifests even without label noise for linear models
(\cite{mei2019generalization}),
and occurs whenever there is \emph{model mis-specification}
between the structure of the true distribution and the model family.
We believe this intuition extends to deep learning as well,
and it is consistent with our experiments.

% Overparameterized models, however, have the capacity to ``memorize the noise'',
% without affecting their global structure as much.
% Moreover, SGD (through unexplained means) manages to find such a minima.
% For overparameterized models, there are many minima which fit the train data,
% an
% and SGD (through unexplained means) is able to ``memorize the noise''
% without destroying the global structure of the model.


\iffalse
\textcolor{blue}{ {\bf Omit this paragraph? --Preetum.}
The situation for linear models is most clear with added label noise.
Here, the fundamental reason why there is high test error
at the interpolation threshold is: when the number of samples $n$
is equal to the dimension $d$ of the linear model,
then there is exactly one model which fits the train set--- and this 
interpolating model is very sensitive
to noise in the train set.
That is, since the linear model is just barely able to fit the train data,
forcing it to fit even slightly-noisy
labels will destroy its global structure, and result in high test error.
However for ``overparameterized'' linear models, when $d \gg n$,
there are many interpolating models that fit the train set,
and selecting the minimum-norm one will often
achieve good test performance\footnote{Note that for linear problems,
the minimum-norm solution is the one found by gradient descent from $0$ initialization.}.
That is, overparameterized linear models are able to
fit the label noise, while maintaining good performance on the distribution.
In general, this situation manifests even without label noise, whenever there is
\emph{model mis-specification} between the structure of the true distribution, and the
linear model family.
}

%% This section follows the linear regression
%%
We believe similar mechanisms may be at work in deep neural networks.
Informally, our intuition is:
for model-sizes at the interpolation threshold, there is effectively
only one model that fits the train data.
This model is very sensitive to [noise in] the specific train samples.
%That is, these models have ``high variance'' in the classical sense.
Overparameterized models, however, have the capacity to ``memorize the noise'',
without affecting their global structure as much.
Moreover, SGD (through unexplained means) manages to find such a minima.
For overparameterized models, there are many minima which fit the train data,
an
and SGD (through unexplained means) is able to ``memorize the noise''
without destroying the global structure of the model.
\fi

% The setting of kernel regression is roughly:
% We draw $n$ samples $(x_i \in \mathcal{X}, y_i \in \R)$
% from a distribution,
% and fix a $d$-dimensional feature map
% $\phi: \mathcal{X} \to \R^d$.
% We want to learn a regressor $f_w(x) := \langle w, \phi(x) \rangle$.
% We pick $\hat{w}$ by minimizing the mean-squared-error over the train samples,
% selecting the minimum norm solution when multiple minima exist.
% That is, $\hat{w} := \phi(X)^\dagger y$
% where $\phi(X)_{i, j} = \phi(x_i)_j$.
% Now, when 
% The test loss of $f_{\hat w}$ 

